\documentclass[journal]{IEEEtran}

\usepackage[T1]{fontenc}
\usepackage{amsmath,amssymb}
\usepackage{cite}
\usepackage{graphicx}
\usepackage[caption=false,font=footnotesize]{subfig}
\usepackage{listings}
\usepackage{url}

\newcommand{\code}[1]{\texttt{#1}}
\newcommand{\pkg}[1]{\texttt{#1}}
\newcommand{\vect}[1]{\boldsymbol{#1}}
\newcommand{\sfive}{\pkg{SymbolicLongMemorySequences.jl}}
\newcommand{\prodacfpath}{Plots/validation/lrd_diagnostics/production/acf}
\newcommand{\prodspectrumpath}{Plots/validation/lrd_diagnostics/production/spectrum}
\newcommand{\marginalpath}{Plots/validation/marginal_control}
\newcommand{\benchpath}{Plots/benchmark/results}
\lstset{
  basicstyle=\ttfamily\footnotesize,
  columns=fullflexible,
  frame=single,
  breaklines=true,
  keepspaces=true,
  showstringspaces=false
} 

\usepackage{xspace}
\newcommand{\ie}{{\em i.e.,}\xspace}
\newcommand{\eg}{{\em e.g.,}\xspace}

\newcommand{\acfsubfig}[2]{%
\subfloat[#2]{%
\IfFileExists{\prodacfpath/#1_symbol_acf.pdf}{%
  \includegraphics[width=0.275\textwidth]{\prodacfpath/#1_symbol_acf.pdf}%
}{%
  \fbox{\parbox{0.255\textwidth}{\centering Missing ACF plot for #2}}%
}}}
\newcommand{\benchfig}[3]{%
\begin{figure}[!t]
\centering
\IfFileExists{\benchpath/#1.pdf}{%
  \includegraphics[width=\columnwidth]{\benchpath/#1.pdf}%
}{%
  \fbox{\parbox{0.92\columnwidth}{\centering Missing generated benchmark plot
  PDF. Run \code{julia --project=benchmark benchmark/benchmarks.jl} in the
  extended benchmark mode and convert the SVG outputs to PDF.}}%
}
\caption{#2}
\label{fig:#3}
\end{figure}}

\title{Reproducible Synthesis of Long-Range Dependent Symbolic Sequences in Julia}

\author{Matthew Roughan, Lewis Mitchell, and Cameron Cornell%
\thanks{The authors are with the School of Mathematical Sciences,
The University of Adelaide, Adelaide, SA 5005, Australia.
E-mail: \{matthew.roughan, lewis.mitchell, cameron.cornell\}@adelaide.edu.au.}}

\begin{document}
\maketitle


\begin{abstract}
Synthetic long-range dependent (LRD) data are useful only when the object being
synthesized matches the object being tested.  This is already difficult for
numerical time series; symbolic data add another source of ambiguity.  A DNA
base, a word, or a security-event class has no natural numerical distance, yet many diagnostics and
models quietly introduce one.  This paper presents \sfive{} (Self-Similar Symbols
Sequence Synthesis),\footnote{LRD and statistical self-similarity are closely
related but not identical: LRD describes persistent dependence, often through a
non-summable covariance tail or low-frequency spectral pole, while statistical
self-similarity describes distributional invariance under aggregation or
rescaling. Many classical models connect them through the Hurst parameter, but
finite symbolic sequences require method-specific diagnostics.} a Julia package
for reproducible synthesis of LRD finite-alphabet sequences.  The package
combines property-based methods, which generate numerical LRD processes and then
symbolize them, with model-based methods, which generate symbolic dependence
directly through finite-state, renewal, excitation, or copy/mutation mechanisms.
The paper takes an operational view: specify the alphabet, marginal frequencies,
local structure when available, nominal LRD mechanism, random seed, and validation
evidence.  We give the mathematical construction of each method,
report one-hot ACF and spectral diagnostics, test marginal control, and benchmark
generation cost.  The result is not a universal symbolic-LRD generator, but a
reproducible workbench for estimator testing, controlled machine-learning stress
tests, privacy-preserving simulation, and studies of symbolic data such as text,
DNA, cyber-event logs, and network measurements.
\end{abstract}

\begin{IEEEkeywords}
Long-range dependence, symbolic sequences, synthetic data, self-similarity,
Julia, reproducible research.
\end{IEEEkeywords}


\section{Introduction}

Reproducible research needs data that can be shared, regenerated, and varied
under controlled conditions.  In many domains the real data are proprietary,
privacy-sensitive, expensive to label, incomplete, or simply unavailable at the
scale needed for controlled experiments.  Synthetic data can provide ground
truth, sensitivity checks, and repeatable comparisons
\cite{ringberg2008_need}, but only when the generator is explicit about what it
controls and what assumptions it inserts.  Network topology and traffic-matrix
synthesis illustrate the same tension: a useful generator must be controllable
enough for experiments, realistic enough not to mislead, and reproducible enough
for other researchers to interrogate \cite{bowden2014_cold,tune2015_spatiotemporal}.

Long-range dependence (LRD), or long memory, makes this need sharper.  LRD
means that observations far apart in time remain statistically coupled.  In the
covariance setting, a stationary process is often described as LRD when
\begin{equation}
  \gamma(k) \sim c_\gamma k^{-\beta}, \qquad 0 < \beta < 1,
\end{equation}
so the covariance tail is not summable.  Equivalently, many familiar models have
a low-frequency spectral pole, with spectral density proportional to
$f^{-(1-\beta)}$ \cite{samorodnitsky2006_long,pipiras2017_long}.  Averages,
counts, spectra, and estimators can therefore converge much more slowly than
they would for short-memory data.  Validation often needs long sequences and
multiple replicates, which is exactly the regime where real traces are hardest
to obtain and share.

The usual LRD story is numerical: river levels, traffic volumes, queue lengths,
or other scalar time series.  The focus here is different.  Many data streams
with plausible long memory are symbolic rather than naturally numerical: text,
DNA, categorical event logs, cyber-security disclosures, packet-loss states,
social-media events, and finite-state workflow traces.  Generating such
sequences supports several tasks:
\begin{itemize}
\item controlled finite-alphabet examples for LRD estimator development and
finite-sample sensitivity studies;
\item test sets for modeling estimators, including entropy and entropy-rate
estimators, mutual-information diagnostics, and string-matching algorithms;
\item stress tests for sequence models on non-language data with language-like
long context; and
\item privacy-preserving synthetic event streams when real traces cannot be
shared.
\end{itemize}
Modeling the generation mechanisms can also sharpen scientific questions: which
features of a sequence are explained by marginals, by local transition rules, by
renewal structure, by latent numerical dependence, or by genuine long context?

Symbolic LRD is not obtained by simply attaching a Hurst parameter to an
alphabet.  A common shortcut is to map symbols to numbers and then use numerical
LRD tools.  That can be useful, but it introduces a metric that the alphabet may
not possess.  Mapping \(\{\mathtt{A},\mathtt{C},\mathtt{G},\mathtt{T}\}\) to
\(\{1,2,3,4\}\), for instance, makes the numerical distance from
\(\mathtt{A}\) to \(\mathtt{T}\) three times the distance from \(\mathtt{A}\) to
\(\mathtt{C}\).  That geometry is an artifact of the encoding, not of the
biology.  Conversely, reducing a symbolic sequence to a few count series can
hide dependence that lives in order, recurrence, or regime changes.

The task is therefore operational.  A useful symbolic-LRD generator must answer
more than ``what is the nominal \(H\)?''  It must specify the alphabet, the
marginal distribution, any short-range structure that should be controlled, the
mechanism that creates long memory, the finite memory or truncation used in
practice, the random seed and provenance, and the diagnostic used to decide
whether the generated sequence behaves as intended.  A method may match a
marginal distribution exactly while producing weak symbolic ACF diagnostics;
another may have a convincing latent numerical LRD process that is attenuated by
quantization; another may generate heavy-tailed recurrence times without a clean
one-hot spectral pole.  These are not minor implementation details.  They are
part of the scientific meaning of the synthetic data.

We present \sfive{}, a Julia package that makes these choices explicit.  The
package brings together property-based methods, which generate numerical LRD
processes and transform them into symbols, and model-based methods, which build
symbolic dependence directly through finite-memory, regime, renewal,
self-exciting, or copy/mutation mechanisms.  The central goal is reproducible
experimentation: explicit constructors, explicit random-number generators, saved
provenance, fast contract tests, opt-in validation campaigns, and benchmark
artifacts.  The package is not a universal symbolic-LRD generator.  It is a
workbench in which mechanisms can be constructed, compared, diagnosed, and
rejected when they do not support the claim being made.

The contributions of this work are:
\begin{itemize}
\item a reproducible generator framework for finite-alphabet symbolic LRD
sequences;
\item a taxonomy separating property-based latent-process transformations from
model-based symbolic mechanisms;
\item new and adapted symbolic synthesis methods, including finite-memory uses
of LAMP-style processes, a scalable dyadic-history LAMP approximation, and
self-exciting symbolic-mass variants motivated by the failure of a naive
normalized Hawkes translation;
\item explicit control contracts for alphabets, marginal distributions, local
Markov structure, nominal LRD parameters, and finite memory limits;
\item a complexity analysis linking generation cost to sequence length, alphabet
size, history depth, and calibration iterations;
\item validation scripts that diagnose both symbolic one-hot processes and, for
property-based generators, the original numerical latent process.
\end{itemize}

Not every proposed model performs well, and this is part of the paper's
argument.  Negative and partially successful constructions are retained because
they show where intuition from numerical LRD breaks down for symbolic data.  The
paper therefore treats validation evidence as part of the method, not as an
afterthought.

\section{Related Work}

The LRD literature is mature, but not monolithic.  Covariance definitions focus
on the tail of \(\gamma(k)\); spectral definitions focus on the low-frequency
pole; aggregation definitions connect to statistical self-similarity; and
point-process definitions often use count variance or recurrence-time tails
\cite{beran1994_statistics,samorodnitsky2006_long,pipiras2017_long}.  The
distinctions are not academic bookkeeping.  A finite-state chain can have strong
local persistence without asymptotic LRD, while a renewal mechanism can have
heavy-tailed counts without giving a simple symbol-level covariance target.
This is why the package is organized by mechanism rather than by a single
advertised exponent.

Numerical LRD synthesis provides the technical starting point.  Spectral fGn
approximations \cite{paxson1997_fast}, exact and approximate Gaussian simulation
\cite{wood1994_simulation,dietrich1997_fast,dieker2004_simulation,coeurjolly2000_simulation},
wavelet and multiscale methods \cite{roughan2000_real}, intermittent maps
\cite{millan2021_on}, mixtures of short-memory components
\cite{sorbye2017_an}, and heavy-tailed on/off or fractal point-process models
\cite{garrett1994_analysis,lowen1995_estimation,ryu1998_point} all offer ways
to make long memory.  They also show the usual engineering tradeoff.  Exact
Gaussian methods give strong second-order guarantees but can be expensive or
fragile for some covariance forms.  Spectral methods are fast and transparent
but approximate.  Renewal and on/off models are interpretable, but finite runs
can be dominated by missing scales or realization-to-realization variability
\cite{roughan1999_the}.  A symbolic package should expose these tradeoffs rather
than hide them behind a constructor named after \(H\).

Symbolic LRD appears most clearly in DNA and text.  Voss's DNA analysis
\cite{voss1992_evolution} is especially important for this paper because it
uses base-specific indicator processes instead of imposing an arbitrary scalar
encoding.  Related DNA work studies nucleotide correlations, spectra, and
copy/mutation mechanisms
\cite{peng1992_long,li1992_long,li1994_understanding,buldyrev1998_analysis}.
Text studies report long correlations, hierarchy, bursty word use, and rare-word
clustering
\cite{schenkel1993_long,ebeling1994_entropy,montemurro2002_long,alvarez2006_hierarchical,altmann2009_beyond,altmann2012_on,tanaka2016_long},
with more recent work connecting long context to neural sequence models
\cite{belletti2018_factorized,belletti2019_quantifying,tay2021_long}.  This
literature is a useful warning: DNA walks, word counts, rare-word arrival
processes, semantic embeddings, and one-hot indicators can all see different
aspects of the same sequence.

Several symbolic models are close to the constructions used here.  Linear
additive Markov processes \cite{kumar2017_linear} and additive Markov-chain
memory-function models
\cite{melnyk2006_memory,mayzelis2006_additive,melnik2014_entropy} motivate
finite-alphabet mechanisms in which weighted history affects the next-symbol
distribution.  Hawkes-process models for word occurrence \cite{ogura2022_modeling}
provide a bursty self-exciting route, though the categorical one-symbol-per-time
setting changes the normalization problem.  Duplication and
expansion-modification models for DNA-like sequences
\cite{li1994_understanding,koroteev2015_duplication,salgado2012_exact} provide
a growth mechanism.  \sfive{} implements these ideas as transparent synthesis
methods rather than fitted models: parameters are construction controls, not
claims of inference from a real corpus.  Some implementations are therefore new
operational variants of older ideas, especially the dyadic-history LAMP
approximation and the MB4b/MB4c self-exciting mass processes.

Finally, symbolic long memory also touches information-theoretic work on texts
and Hilberg-like laws \cite{debowski2015_hilberg}.  That connection is not the
main focus of this paper, but it reinforces the same point: for symbolic data,
the object of interest is often not one scalar time series but a collection of
recurrences, counts, transitions, and transformations.  A useful synthesis tool
must make those choices visible.

\section{Operational Design}

The introduction frames symbolic synthesis as an operational specification
problem.  This section turns that frame into the package design: what must a user
define before a synthetic symbolic sequence is useful for an experiment?  Length
and alphabet are not enough.  The user must decide the
observable symbol space, the intended long-memory mechanism, the marginal
frequencies, any short-range structure that should survive generation, and the
random source needed for reproducibility.  We write a generator schematically as
\begin{equation}
  g =
  (\mathcal{A}, n, \theta_{\mathrm{LRD}}, \vect{p}, \mathcal{L}, d,
   \code{rng}, \text{method}),
\end{equation}
where \(\mathcal{A}\) is the ordered alphabet, \(n\) is the sequence length,
\(\theta_{\mathrm{LRD}}\) is a method-specific long-memory parameter,
\(\vect{p}\) is an optional marginal distribution, \(\mathcal{L}\) is optional
local structure such as Markov transition matrices, \(d\) is a configured or
effective memory cutoff when applicable, and \code{rng} records the reproducible
random source.  Not every method supports every control.  The API therefore
records supported controls rather than presenting all methods as interchangeable.
Throughout the paper, bold symbols in math mode denote vectors; scalar
components of those vectors are written with subscripts.

\subsection{LRD Parameters and Self-Similarity}

For a second-order stationary numerical process, the common covariance
parameterization is
\begin{equation}
  \gamma(h) \sim c_\gamma h^{-\beta}, \qquad 0<\beta<1,
\end{equation}
where \(\gamma(h)=\operatorname{Cov}(X_t,X_{t+h})\).\footnote{We use
\(\sim\) in its standard asymptotic-equivalence sense: for positive functions
\(f\) and \(g\), \(f(k)\sim g(k)\) as \(k\to\infty\) means
\(\lim_{k\to\infty} f(k)/g(k)=1\). Thus a relation such as
\(\gamma(k)\sim c_\gamma k^{-\beta}\) is a large-lag tail statement including
its multiplicative constant, not a claim that finite validation plots exactly
match a power law at every displayed lag.}  The corresponding
low-frequency spectral form is
\begin{equation}
  f(\lambda) \sim c_f |\lambda|^{-\alpha_s},
  \qquad \alpha_s = 1-\beta,
\end{equation}
and the usual Hurst-parameter relationship is
\begin{equation}
  H = 1-\frac{\beta}{2}
    = \frac{1+\alpha_s}{2},
  \qquad 1/2 < H < 1.
\end{equation}
Equivalently, \(\beta=2-2H\) and \(\alpha_s=2H-1\).  These equations are a
translation guide among common numerical conventions.  They are not a guarantee
that a finite symbolic generator with a nominal \(H\) will display the same
slope in every diagnostic.

The notation needs care because the literature reuses letters.  In this paper
\(\alpha_s\) denotes a spectral exponent, while some generator descriptions use
\(\alpha\) for a heavy-tail exponent, for example the Pareto sojourn exponent in
MB2.  For ideal heavy-tailed on/off constructions with \(1<\alpha<2\), the
aggregate count-process heuristic is \(H=(3-\alpha)/2\), which is not the same
relationship as \(\alpha_s=2H-1\).  The package documentation therefore explains
parameters method by method rather than treating \(H\), \(\alpha\), and \(\beta\)
as interchangeable labels.

LRD is closely connected to statistical self-similarity.  If a process is
aggregated into blocks of size \(m\), self-similar scaling means that suitably
normalized block averages retain the same large-scale distributional form.  For
symbolic sequences the observable object may be a symbol-indicator process,
symbol counts in blocks, recurrence times, or a latent numerical driver.  For this reason \sfive{} reports mechanism, transformation, and diagnostic
together: a symbol sequence can be self-similar in one operational view and
almost invisible in another.

\subsection{Marginal and Local Controls}

The marginal distribution \(p_j=\Pr(X_t=a_j)\) controls how often each symbol
appears.  It matters before we even discuss LRD.  First, it determines the
variance of one-hot diagnostics: rare symbols give noisy autocorrelations.
Second, it is often a domain constraint.  DNA is the clean example.  Chargaff's
rules concern the relative frequencies of A, C, G, and T, and real genomes have
species-specific base compositions as well as differences between species
\cite{chargaff1950_chemical}.  A symbolic DNA generator that cannot control at
least the marginal distribution is missing a basic biological feature.  Text has
an equally familiar version of the same issue: word frequencies often follow a
Zipf-like distribution, so a synthetic text-like sequence with uniform word
frequencies would miss a first-order property of the domain
\cite{zipf1949_human}.  Similar concerns appear as event-frequency imbalance in
logs.  Third, uncontrolled marginals can make a synthetic benchmark too easy: a
model may learn the symbol imbalance rather than the long memory.

Short-range structure is a different knob.  A first-order Markov specification
uses a transition matrix
\begin{equation}
  P_{ij}=\Pr(X_t=a_j\mid X_{t-1}=a_i)
\end{equation}
to prescribe local repetition, alternation, or symbol-to-symbol preference.  In
\sfive{} this appears directly in methods such as PB3 and MB2, and indirectly
through transition or excitation matrices in additive and self-exciting models.
Controlling \(P\) is not the same as controlling LRD: a finite-state Markov
chain with a fixed transition matrix is short-memory under ordinary conditions.
The long-memory mechanism must enter through a persistent latent driver, a
heavy-tailed regime duration, a power-law history kernel, or another scale-free
construction.

The design goal follows directly from the paper's narrative: simultaneous
control with explicit limits.  A useful symbolic LRD generator should let the user say: use this alphabet; make these symbols
common or rare; preserve this local Markov behavior when the method supports it; and inject
long memory through this documented mechanism.  These controls can
conflict.  Strong marginal calibration can attenuate latent dependence; Markov
regimes with identical stationary distributions can hide LRD from one-hot
diagnostics; finite memory \(d\) can truncate a nominal power law.  The package
therefore reports both capability and caveat in constructors, documentation,
validation plots, and provenance metadata.

\subsection{Interface and Evidence}

The public contract is intentionally small:
\begin{equation}
  \code{generate(g, n; rng)} \rightarrow \vect{X}_{1:n}.
\end{equation}
Generators preserve alphabet element types and reject invalid probabilities,
transition matrices, and parameter ranges at construction time.  Convenience
functions construct standard cases by method identifier, while method-specific
constructors remain available when the scientific details matter.  The short API
supports routine experiments; explicit constructors keep the method assumptions
visible.

Property-based generators also implement
\begin{equation}
  \code{generate\_with\_latent}(g,n;\code{rng}),
\end{equation}
which returns both the emitted symbols and the latent numerical matrix.  The
latent output is needed for validation.  If the symbolic diagnostics look weak,
it helps distinguish a weak numerical source from dependence lost in the
transformation or from finite-sample variability.

The reproducibility story also requires separating correctness evidence from
scientific evidence.  Unit tests
cover construction, reproducibility, type preservation, and finite-sample
contracts.  Validation scripts produce longer Monte Carlo diagnostics: one-hot
autocorrelation, one-hot spectra, marginal control, local transition control,
and comparison conventions compatible with LongMemory.jl.  Benchmarks live in a
separate project environment and run only on explicit request.  This separation keeps ordinary package tests short while retaining the
longer validation evidence used in this paper.

\section{Methods}

The methods are grouped by how they create dependence.  Property-based methods
start with a numerical LRD process and then ask how much dependence survives the
symbolizer.  Model-based methods generate symbols directly, so their claims are
tied to the symbolic mechanism itself.  Table~\ref{tab:methods} summarizes the
implemented methods, their main controls, computational cost, and the caveat a
user should keep in mind before trusting a plot.

\begin{table*}[!t]
\caption{Implemented symbolic LRD synthesis methods in \sfive{}. Here $n$ is the
sequence length, $K$ is the alphabet size, $d$ is a configured history depth,
and $I$ is the number of PB2 calibration iterations. Costs describe generation
after construction for the standard implementations.}
\label{tab:methods}
\centering
\footnotesize
\begin{tabular}{lp{0.18\textwidth}p{0.18\textwidth}lp{0.26\textwidth}}
\hline
ID & Mechanism & Main controls & Typical cost & Main caveat \\
\hline
PB1 & Spectral fGn + rank quantization & $H$, marginal & $O(n\log n)$ & local structure induced by quantization \\
PB2 & Latent Gaussian categorical argmax & $H$, marginal, $I$ & $O(nKI)$ & argmax transform is nonlinear and calibrated empirically \\
PB3 & LRD regime driver + Markov chain & $H$, regime weights, $P^{(r)}$ & $O(n\log n+nK)$ & one-hot diagnostics may miss hidden regimes \\
PB4 & Intermittent map + rank quantization & map exponent, marginal & $O(n\log n)$ & finite laminar scales drive behavior \\
\hline
MB1a & Exact finite-history additive row mixture & $\beta$, $d$, $P$, marginal & $O(nd)$ & finite memory truncates asymptotic LRD \\
MB1b & Dyadic-bucket additive row mixture & $\beta$, $d$, $P$, marginal & $O(nK\log d)$ & dyadic approximation changes exact history weights \\
MB1c & Centered additive memory function & $\beta$, $d$, $\rho$, marginal & $O(nd)$ & clipping/renormalization affects strong memory \\
MB2 & Heavy-tailed regime sojourns & $\alpha$, $P^{(r)}$, $Q$ & $O(nK)$ & count-process theory is only a guide for symbols \\
MB3 & Competing fractal renewal streams & tail exponent, rates & $O(nK)$ & missing-scale effects can be severe \\
MB4a & Normalized Hawkes-style excitation & $\beta$, $d$, $\vect{\mu}$, $A$ & $O(nKd)$ & normalized kernel can look short-memory \\
MB4b & Self-exciting symbol mass & $\beta$, $d$, $\vect{p}$, $b$, $\eta$, $s$, $A$ & $O(nKd)$ & strong excitation can distort marginals \\
MB4c & Log-contrast self-exciting mass & $\beta$, $d$, $\vect{p}$, $b$, $\eta$, $s$, $a$, $A$ & $O(nKd)$ or $O(n^2K)$ & contrast can cause symbol lock-in \\
MB5 & Power-law lag copy/mutation growth & $\alpha$, $d_{\max}$, mutation & $O(n\log d)$ & growth mechanism can overproduce local patches \\
\hline
\end{tabular}
\end{table*}

\subsection{Property-Based Methods}

Property-based methods are attractive because numerical LRD synthesis is well
understood.  They are also dangerous in a useful way: the transformation can
change the dependence we thought we had.  The package therefore treats the
latent process and the symbolizer as separate objects.

Let $\mathcal{A}=\{a_1,\ldots,a_K\}$ be the ordered alphabet and let
\(\vect{p}=(p_1,\ldots,p_K)\) be a target marginal distribution where applicable.
Property-based generators first create a real-valued latent matrix
$Z\in\mathbb{R}^{W\times n}$ and then apply a symbolizer
$T:\mathbb{R}^{W\times n}\rightarrow\mathcal{A}^n$. The implementation exposes
this split through
\[
  (\vect{X}_{1:n},Z) = \code{generate\_with\_latent}(g,n;\code{rng}),
\]
so validation can ask whether the LRD signal is present in $Z$, lost in $T$, or
only weakly visible after finite-sample averaging.

\paragraph{PB1: spectral fGn with quantization.}
PB1 generates a single latent series $Z_t$ using Paxson's approximate spectral
fractional Gaussian noise method \cite{paxson1997_fast}. On the Fourier grid,
the target spectral shape is
\begin{equation}
  S(f_j) \propto |f_j|^{1-2H}, \qquad H\in(1/2,1),
\end{equation}
with random phases and Hermitian symmetry giving a real inverse FFT. The
implementation normalizes the latent series to zero sample mean and unit sample
variance. The symbolizer is deterministic rank binning: choose integer bin
counts $c_i$ with $\sum_i c_i=n$ and $c_i\approx n p_i$, sort the latent values,
and assign the smallest $c_1$ values to $a_1$, the next $c_2$ values to $a_2$,
and so on. Thus PB1 has strong finite-sample marginal control, but short-range
symbol structure is induced rather than directly prescribed.

\paragraph{PB2: latent Gaussian categorical model.}
PB2 uses $K$ independent fGn streams $Z_{i,t}$ with common Hurst parameter $H$.
The emitted symbol is
\begin{equation}
  X_t = a_{\arg\max_i \{Z_{i,t}+\theta_i\}},
\end{equation}
where offsets \(\theta_i\) are entries of the offset vector \(\vect{\theta}\),
which is calibrated on the realized latent sample to
approximate the target marginal distribution. \sfive{} uses the iterative update
\begin{equation}
  \theta_i \leftarrow \theta_i
      + \eta\left[\log(p_i+\epsilon)-\log(\hat p_i+\epsilon)\right],
\end{equation}
followed by centering the offsets, where \(\hat p_i\) is the current empirical
choice probability and \(\hat{\vect{p}}\) is the empirical choice-probability
vector. This is a practical finite-sample adaptation of latent
Gaussian categorical modeling \cite{gal2015_latent}; it does not solve the
exact multivariate Gaussian choice probabilities.

\paragraph{PB3: latent driver with Markov regimes.}
PB3 creates a one-dimensional long-memory driver $Z_t$ using either the spectral
source above or a Haar-like cascade. The driver is rank-binned into regimes
\(R_t\in\{1,\ldots,M\}\) with target regime-weight vector \(\vect{\pi}\).
Conditional on the
current regime, symbols follow a first-order Markov chain:
\begin{equation}
  \Pr(X_t=a_j\mid X_{t-1}=a_i,R_t=r) = P^{(r)}_{ij}.
\end{equation}
The local bigram mechanism is therefore explicit through the matrices
$P^{(r)}$, while large-scale persistence is inherited from the regime driver.
When regimes have nearly identical stationary distributions, one-hot symbol
diagnostics may hide the latent LRD; validation therefore reports the latent
driver as well as the emitted symbols.

\paragraph{PB4: intermittent-map quantization.}
PB4 uses a Pomeau-Manneville-style intermittent map
\begin{equation}
  Z_{t+1} = (Z_t + Z_t^z) \bmod 1, \qquad z>1,
\end{equation}
after a burn-in period. The long laminar episodes near zero create broad
finite-scale dependence in the latent driver. \sfive{} then applies the same
rank-binning symbolizer as PB1. This method is motivated by intermittent-map
approaches to genomic-sequence statistics \cite{provata2012_coupled}, but the
package treats it as a latent finite-sample mechanism rather than an exact
closed-form symbolic LRD construction.

The common risk is that nonlinear transformations can change apparent
long-memory behavior. The package therefore validates both the latent numerical
process and the final one-hot symbolic process.

\subsection{Model-Based Methods}

Model-based methods generate symbols directly.  They avoid the arbitrary
numeric-encoding problem, but they must earn their LRD behavior through the
symbolic mechanism itself.  This makes their caveats more visible: finite memory,
renewal overdispersion, normalization, and copy/mutation patches all show up in
validation.

\paragraph{MB1a: finite-history LAMP.}
MB1a implements a finite-history linear-additive Markov process
\cite{kumar2017_linear}. Given transition matrix $P$, target marginal \(\vect{p}\),
history depth $d$, and normalized weights $w_\ell$, the next-symbol
distribution vector \(\vect{q}_t\) has entries
\begin{equation}
  q_t(j) = (1-\epsilon)\sum_{\ell=1}^{m_t} w_{\ell,t}
             P_{I(X_{t-\ell}),j}
          + \epsilon p_j + b_{t,j},
\end{equation}
where $m_t=\min(d,t-1)$, $I(X)$ maps symbols to indices, and $b_{t,j}$ accounts
for missing pre-history mass by assigning it according to \(\vect{p}\). \sfive{} uses
power-law weights over observed lags,
\begin{equation}
  w_\ell \propto \ell^{-(1+\beta)}
\end{equation}
for the LAMP-style row mixture. The small innovation parameter $\epsilon$
prevents finite-history absorption and preserves a direct marginal-control
path. MB1a is exact for the configured finite memory, but it is not a proof that
the finite sequence has asymptotic LRD beyond the memory window.

\paragraph{MB1b: dyadic-bucket LAMP.}
MB1b approximates the MB1a history sum with dyadic lag buckets
\begin{equation}
  B_b = \{2^b,\ldots,2^{b+1}-1\}.
\end{equation}
For each bucket, the generator maintains a count vector \(\vect{h}_b\) of
historical symbols
and a total bucket weight $W_b=\sum_{\ell\in B_b} w_\ell$. The row-mixture term
is approximated by
\begin{equation}
  \sum_b W_b \sum_i \hat h_{b,i} P_{ij},
\end{equation}
where $\hat h_{b,i}$ is the empirical symbol mix in bucket $b$. This preserves
the power-law age profile at dyadic resolution while avoiding an explicit scan
over all $d$ lags. It is a scalable approximation, not an identical
finite-history LAMP.

\paragraph{MB1c: centered additive Markov chain.}
MB1c follows additive Markov-chain memory-function work
\cite{melnyk2006_memory,mayzelis2006_additive,melnik2014_entropy}. It writes
the next-symbol probability vector as a perturbation around the target marginal:
\begin{equation}
  q_t(j) = p_j + \rho\sum_{\ell=1}^{m_t} v_{\ell,t}
          \left(\mathbf{1}\{X_{t-\ell}=a_j\}-p_j\right),
\end{equation}
with $v_\ell\propto \ell^{-\beta}$ and strength $\rho\in[0,1]$. The centered
form makes the intended recurrence effect clearer than a row mixture: observing
$a_j$ in the past increases the future probability of $a_j$ relative to its
baseline. Probabilities are clipped and renormalized when finite histories and
large $\rho$ would otherwise leave the simplex.

\paragraph{MB2: heavy-tailed on/off Markov chain.}
MB2 has a regime process $R_t$ with heavy-tailed sojourns. A sojourn length $L$
has Pareto-like tail
\begin{equation}
  \Pr(L>x) = (L_{\min}/x)^\alpha,\qquad 1<\alpha<2,
\end{equation}
above the scale $L_{\min}$. At regime boundaries, a transition matrix $Q$
selects the next regime. Within a regime,
\begin{equation}
  \Pr(X_t=a_j\mid X_{t-1}=a_i,R_t=r)=P^{(r)}_{ij}.
\end{equation}
Heavy-tailed regime durations create persistent occupancy, while the
within-regime matrices control local symbol dynamics. The nominal count-process
relationship is $H=(3-\alpha)/2$ for ideal on/off constructions
\cite{garrett1994_analysis,ryu1998_point}.

\paragraph{MB3: fractal symbol sequence.}
MB3 assigns each symbol $a_i$ an independent event stream with heavy-tailed
inter-arrival times. If \(T_i\) is the entry of the pending-time vector
\(\vect{T}\) for symbol \(i\), the
emitted symbol at sequence step $t$ is the one with the smallest pending event
time:
\begin{equation}
  X_t = a_{\arg\min_i T_i}.
\end{equation}
After emitting $a_i$, its pending time is advanced by another heavy-tailed
inter-arrival sample. This adapts fractal renewal and shot-noise ideas to a
finite alphabet \cite{lowen1995_estimation,ryu1998_point}. The known
missing-scales problem for naive fractal renewals \cite{roughan1999_the} is a
reason to validate the realized scale range rather than rely only on the tail
parameter. MB3's rate vector is therefore an asymptotic intensity control: in a
finite sequence, heavy-tailed renewal-count fluctuations can dominate the
nominal rate ratios.

\paragraph{MB4a, MB4b, and MB4c: discrete-time self-exciting symbols.}
The first MB4 construction, now denoted MB4a and implemented as
\code{HawkesSymbol}, is a finite-history symbolic analogue of Hawkes
word-occurrence models \cite{ogura2022_modeling}. It follows the Hawkes idea of
an immigration or baseline component plus self-excitation from past events. At
time $t$ each symbol has an intensity
\begin{equation}
  \lambda_{t,j} =
  \mu_j + \sum_{\ell=1}^{m_t} \bar{g}(\ell)
          A_{I(X_{t-\ell}),j},
  \qquad
  \bar{g}(\ell)=\frac{(\ell+c)^{-\beta}}
                     {\sum_{r=1}^{d}(r+c)^{-\beta}},
\end{equation}
where \(m_t=\min(d,t-1)\), \(\vect{\lambda}_t\) is the intensity vector,
\(\vect{\mu}\) is a positive baseline vector, and $A$ is a nonnegative excitation
matrix. The emitted symbol is sampled from normalized intensities:
\begin{equation}
  \Pr(X_t=a_j\mid X_{1:t-1}) =
  \lambda_{t,j}/\sum_r\lambda_{t,r}.
\end{equation}
Identity-like $A$ creates bursty repetition; off-diagonal entries create
cross-symbol triggering. This is close to a discrete-time marked or multivariate
Hawkes intuition, but it differs from a continuous-time Hawkes process because
exactly one categorical symbol is emitted at every time step.

Validation exposed a weakness in MB4a. Since \(\bar{g}\) has unit mass, the
history component has bounded total scale. In the uniform identity-excitation
case a first-order linearization has the form
\begin{equation}
  Z_t \approx a \sum_{\ell=1}^{d}\bar{g}(\ell)Z_{t-\ell}+\varepsilon_t,
  \qquad 0<a<1,
\end{equation}
where the baseline and categorical normalization keep $a$ below one. The
low-frequency spectrum can then have a plateau rather than a visible power-law
pole, so centered one-hot diagnostics can look almost white even when short-run
burstiness is present. Increasing the identity-excitation amplitude does not by
itself remove this normalization effect.

MB4b, implemented as \code{SelfExcitingMass}, keeps the discrete-time symbolic
setting but changes what is normalized. Let \(\vect{p}\) be the target/default
marginal, let $b\ge 0$ be a default mass, let \(\eta\ge 0\) be an excitation
strength, and let $s\ge 0$ be a smoothing floor. MB4b forms unnormalized masses
\begin{equation}
  M_{t,j} = b p_j + \eta\left(s p_j +
  \sum_{\ell=1}^{m_t} (\ell+c)^{-\beta} A_{I(X_{t-\ell}),j}\right),
\end{equation}
and only then samples
\begin{equation}
  \Pr(X_t=a_j\mid X_{1:t-1}) = M_{t,j}/\sum_r M_{t,r}.
\end{equation}
The key difference is that the power-law kernel is not normalized in the memory
mass. For \(0<\beta<1\), the represented memory mass grows like
\(d^{1-\beta}\) over the finite range, allowing self-excitation to dominate the
small default component once enough history exists. The default and smoothing
terms are retained to make the startup distribution normalizable and to prevent
symbols from becoming permanently unreachable. MB4b is therefore best viewed as
our discrete-time categorical adaptation of Hawkes-style self-excitation, rather
than a standard Hawkes model already present in the literature. The closest
literature precedent is MB4a's Hawkes/word-recurrence motivation; MB4b is the
package modification introduced because the traditional normalized categorical
translation gave weak LRD diagnostics.

The remaining weakness of MB4b is that the final categorical normalization still
uses the total mass \(\sum_r M_{t,r}\). When the identity-excitation model has
nearly uniform symbol usage, much of the accumulated long-memory mass is common
mode across symbols. MB4c, implemented as \code{LogitSelfExcitingMass}, keeps the
same unnormalized mass \(M_{t,j}\) but converts it to a centered log contrast
before sampling. Let \(a\ge 0\) be a contrast strength and \(\epsilon>0\) a small
logarithmic floor. Define
\begin{equation}
  \theta_{t,j} = a\left[\log(M_{t,j}+\epsilon)
  - \sum_r p_r\log(M_{t,r}+\epsilon)\right].
\end{equation}
MB4c then samples according to
\begin{equation}
  \Pr(X_t=a_j\mid X_{1:t-1}) =
  \frac{p_j\exp(\theta_{t,j})}{\sum_r p_r\exp(\theta_{t,r})}.
\end{equation}
The subtraction removes common-mode growth in the history mass before the final
probability normalization. Conceptually, the power-law kernel need not have a
finite memory limit; the implementation parameter \(d\) is a computational
truncation, and \code{d = nothing} in MB4c uses all previous generated symbols.
The offset \(c\) is likewise not a second LRD parameter: it regularizes short
lags through \((\ell+c)^{-\beta}\), while \(\beta\) controls the asymptotic
power-law decay. Preliminary validation shows the expected tradeoff. Increasing
\(a\) can steepen the low-frequency spectrum, but too large a contrast creates
symbol lock-in and poor finite-sample marginals. We therefore treat MB4c as a
new experimental variant and use it to expose the remaining design problem rather
than as a final closed-form solution.

\paragraph{MB5: duplication-mutation growth.}
MB5 starts from an iid seed sequence and grows by copying a previous symbol at
a random lag. The lag $D_t$ is drawn from a truncated power law,
\begin{equation}
  \Pr(D_t=d) \propto d^{-\alpha},
  \qquad 1\le d\le d_{\max}.
\end{equation}
With probability $1-\mu$, the next symbol is copied,
$X_t=X_{t-D_t}$; with probability $\mu$, it is mutated by drawing from the
target marginal or an alternative symbol distribution. This is motivated by
duplication-mutation and expansion-modification ideas in DNA-like symbolic
sequences \cite{li1994_understanding,koroteev2015_duplication,salgado2012_exact}.
\sfive{} uses the power-law lag-copy version because validation showed that a
block-length-only version produced local patches without the intended symbolic
autocorrelation decay.

\section{Validation Results}
\label{sec:validation}

\subsection{ACF and Spectra for Symbolic Sequences}

The introduction argues that symbolic LRD claims depend on the diagnostic used.
The first validation question is therefore how to compute an ACF for objects that
are not numbers.  A symbolic sequence has no canonical numerical scale.  For example,
replacing
\(\{\mathtt{A},\mathtt{C},\mathtt{G},\mathtt{T}\}\) by \(\{1,2,3,4\}\) would make
the distance between \(\mathtt{A}\) and \(\mathtt{T}\) three times the distance
between \(\mathtt{A}\) and \(\mathtt{C}\), even though no such distance is part
of the alphabet. \sfive{} therefore validates symbolic dependence through centered
indicator processes, following the spirit of Voss's DNA indicator method
\cite{voss1992_evolution}.

Let \(\vect{X}_{1:n}=(X_1,\ldots,X_n)\in\mathcal{A}^n\), where
\(\mathcal{A}=\{a_1,\ldots,a_K\}\). For each symbol \(a_j\), define the binary
indicator and its centered version
\begin{equation}
  B_t(a_j)=\mathbf{1}\{X_t=a_j\}, \qquad
  Y_t(a_j)=B_t(a_j)-\bar B(a_j),
\end{equation}
where \(\bar B(a_j)=n^{-1}\sum_{t=1}^n B_t(a_j)\). The sample autocovariance for
symbol \(a_j\) at lag \(h\) is then
\begin{equation}
  \hat\gamma_j(h)=
  \frac{1}{n-h}\sum_{t=1}^{n-h}Y_t(a_j)Y_{t+h}(a_j),
  \qquad h=0,\ldots,n-1,
\end{equation}
and the corresponding autocorrelation is
\begin{equation}
  \hat\rho_j(h)=\hat\gamma_j(h)/\hat\gamma_j(0),
\end{equation}
when \(\hat\gamma_j(0)>0\). Symbols with zero empirical variance are omitted.
The reported symbolic ACF is an average across the remaining symbol indicators,
usually after retaining positive values for log-log display:
\begin{equation}
  \hat\rho_{\mathrm{sym}}(h)=
  \frac{1}{|\mathcal{J}|}\sum_{j\in\mathcal{J}}\hat\rho_j(h),
  \qquad
  \mathcal{J}=\{j:\hat\gamma_j(0)>0\}.
\end{equation}
This is not the only possible symbolic diagnostic, but it has a clear
interpretation: it asks whether occurrences of the same symbol remain
correlated over long lags.

The spectral diagnostic uses the same centered indicator series. For Fourier
frequencies \(f_m=m/n\), a simple periodogram for symbol \(a_j\) is
\begin{equation}
  \hat S_j(f_m)=
  \frac{1}{n}\left|
    \sum_{t=1}^n Y_t(a_j)\exp(-2\pi i f_m t)
  \right|^2 .
\end{equation}
\sfive{} averages these symbol-wise periodograms across \(j\in\mathcal{J}\), then
plots the result on log-log axes. A symbolic sequence with covariance-like LRD
in its indicator processes should show slowly decaying positive ACF values and
elevated low-frequency spectral power, subject to finite-sample noise.

As a toy example, take
\[
  \vect{X}_{1:8}=(a,b,a,a,c,b,a,c).
\]
For symbol \(a\),
\[
  \vect{B}(a)=(1,0,1,1,0,0,1,0), \qquad \bar B(a)=4/8,
\]
so
\[
  \vect{Y}(a)=(1/2,-1/2,1/2,1/2,-1/2,-1/2,1/2,-1/2).
\]
The lag-one autocovariance contribution is the average of
\[
  \vect{c}_1=(-1/4,-1/4,1/4,-1/4,1/4,-1/4,-1/4),
\]
which is negative for this short sequence; the symbol \(a\) tends not to repeat
at lag one except for the pair at positions 3 and 4. The same calculation is
performed for \(b\) and \(c\), and the symbol-wise ACFs are averaged. For a real
validation run this procedure is applied to much longer sequences and many
lags; the toy example is only meant to show the transformation from symbols to
numerical diagnostics.

The plotted dashed reference lines show nominal power-law behavior and
finite-sample interpretation limits.  They are visual guides, not estimator
guarantees.  Centered one-hot diagnostics also have limitations: they can miss
dependence that appears only in higher-order patterns, changing transition
matrices, or recurrence-time distributions.  This limitation reflects the problem rather than the implementation: symbolic
LRD cannot be reduced to one statistic, so the diagnostic must be stated beside
the generated data.

\subsection{Marginal-Control Validation}

Because the paper treats synthetic data as a controlled experimental object,
LRD diagnostics are only meaningful if the generator also respects the requested
symbol frequencies well enough for the intended experiment. \sfive{} therefore
includes a separate marginal-control validation. The paper-facing case uses a
uniform categorical target with \(K=8\), generates \(n=100000\) symbols for each
replicate, and drops the first and last 10\% before computing frequencies. The
trimmed window reduces sensitivity to initialization and terminal edge effects,
leaving 80000 observations per replicate in the reported run.

The uniform target has two advantages. First, if a generator can produce a
fine uniform categorical distribution, other marginals can be approximated by
aggregating fine bins; direct nonuniform targets are still tested elsewhere in
the validation script, but the uniform case gives the most comparable shared
diagnostic.
Second, the uniform case avoids small expected cell counts. Classical
chi-squared frequency tests become unreliable when bins have low expected
counts, a common problem for skewed symbolic distributions.

For each replicate, the validation computes
\begin{equation}
  \chi^2_{\mathrm{obs}} = \sum_{j=1}^{K}\frac{(N_j - m/K)^2}{m/K},
\end{equation}
where \(m\) is the trimmed sample size and \(N_j\) is the count for symbol
\(a_j\). The reported p-values use the iid multinomial chi-squared reference
with \(K-1\) degrees of freedom. This reference is conservative as
an interpretation aid: the generated sequences are dependent, so the p-values
are frequency diagnostics rather than exact hypothesis-test probabilities. The
issue is especially important for LRD generators because empirical averages can
converge more slowly than the iid multinomial model assumes. Thus a strict
chi-squared test may reject too often even when the generator has the intended
long-run marginal, although large values remain indicative of a control problem.

The validation table also reports an approximate effective-sample-size (ESS)
correction. For symbol \(a_j\), define the centered indicator
\(Y_{j,t}=\mathbb{1}\{X_t=a_j\}-\hat p_j\) on the trimmed sequence. Let
\(\hat\rho_j(h)\) be its sample autocorrelation at lag \(h\). We use an
initial-positive truncation. That is, \(h_j^\star\) is the largest lag before
the first non-positive value in the sequence
\(\hat\rho_j(1),\hat\rho_j(2),\ldots\), and
\begin{equation}
  \hat\tau_j = 1 + 2\sum_{h=1}^{h_j^\star}\hat\rho_j(h),
\end{equation}
with \(\hat\tau_j\ge 1\), and define
\begin{equation}
  \hat m_{\mathrm{eff}} =
  \min_{1\le j\le K}\frac{m}{\hat\tau_j}.
\end{equation}
The ESS-adjusted statistic is then
\begin{equation}
  \chi^2_{\mathrm{ESS}} = \chi^2_{\mathrm{obs}}\frac{\hat m_{\mathrm{eff}}}{m}.
\end{equation}
This scalar correction is not a full dependent multinomial theory. It uses the
smallest symbol-level ESS, so it is conservative, but it helps identify cases
where an apparently large frequency deviation is partly explained by long
dependence in the symbol indicators.
For MB1a and MB1b, the marginal-control validation uses the package's
repeat-biased identity/dyad transition helper, with repeat probability 0.4.
This keeps a history-repetition mechanism while making \(\vect{p}\) the
stationary distribution of the transition matrix. A pure identity matrix is
still a useful stress case, but it can preserve early finite-sample imbalances
for a long time.
The retained table also reports full-sequence marginal error, because some
methods calibrate or rank-bin the entire generated sample while the trimmed
interior window is allowed to fluctuate.

A better formal test would calibrate the null distribution under dependence.
Possible extensions include block or subsampling tests, or a parametric Monte
Carlo envelope generated from the same configured generator. These extensions
are left for later work. The simple chi-squared reference and the ESS-adjusted
diagnostic remain useful because they are transparent,
comparable across methods, and sensitive to obvious marginal-control failures.

\begin{table*}[!t]
\caption{Uniform marginal validation for \(K=8\), \(n=100000\), 20 replicates,
and a trimmed window of 80000 observations per replicate. The iid p-values use
the multinomial chi-squared reference. The ESS p-values scale the statistic by
the conservative effective sample size estimated from one-hot symbol
autocorrelations. Both should be read as diagnostics under dependence. TV
denotes total variation distance from the target marginal.}
\label{tab:uniform-marginal}
\centering
\footnotesize
\begin{tabular}{lrrrrr}
\hline
Method & Mean \(\chi^2_{\mathrm{obs}}\) & Median iid \(p\) & Median ESS \(p\) & Full TV & Trim TV \\
\hline
\code{PB1} \code{SpectralFGN} & 39.65 & 0.040 & 0.997 & 0.0000 & 0.0074 \\
\code{PB2} \code{LGCM} & 87.69 & \(4.55{\times}10^{-14}\) & 0.966 & 0.0001 & 0.0135 \\
\code{PB3} \code{WaveletMarkov} & 6.84 & 0.514 & 0.525 & 0.0034 & 0.0036 \\
\code{PB4} \code{IntermittentMapSymbols} & 25.18 & 0.065 & 1.000 & 0.0000 & 0.0053 \\
\code{MB1a} \code{LAMP} & 18.92 & 0.021 & 0.445 & 0.0053 & 0.0060 \\
\code{MB1b} \code{DyadicLAMP} & 15.52 & 0.035 & 0.486 & 0.0048 & 0.0055 \\
\code{MB1c} \code{CalibratedAdditiveMarkov} & 7.43 & 0.493 & 0.508 & 0.0035 & 0.0040 \\
\code{MB2} \code{OnOffMarkov} & 6.33 & 0.567 & 0.583 & 0.0032 & 0.0035 \\
\code{MB3} \code{FSS} & 203.01 & \(5.48{\times}10^{-17}\) & \(5.48{\times}10^{-17}\) & 0.0160 & 0.0175 \\
\code{MB4a} \code{HawkesSymbol} & 7.43 & 0.493 & 0.508 & 0.0035 & 0.0040 \\
\code{MB4b} \code{SelfExcitingMass} & 7.43 & 0.493 & 0.508 & 0.0035 & 0.0040 \\
\code{MB4c} \code{LogitSelfExcitingMass} & 7.43 & 0.493 & 0.508 & 0.0035 & 0.0040 \\
\code{MB5} \code{DuplicationMutation} & 6.81 & 0.531 & 0.545 & 0.0033 & 0.0035 \\
\hline
\end{tabular}
\end{table*}

The table shows three different phenomena. PB1 and PB4 use rank binning and
therefore match the full sample exactly in this uniform case, but the interior
trimmed window can wander because the sequence is dependent. PB2 similarly
calibrates the whole generated sample very closely, while the trimmed window has
larger deviations. Methods with explicit iid or stationary uniform mechanisms,
including PB3, MB1a, MB1b, MB1c, MB2, MB4a, MB4b, MB4c, and MB5 in the tested configurations,
mostly match the marginal closely in total-variation terms. MB3 is the clear
exception. Its independent heavy-tailed renewal streams show the largest
replicate variability in both the full and trimmed samples, and the ESS
correction does not change the conclusion: the estimated one-hot ESS is the full
trimmed length in this run. This indicates that the failure is not primarily a
short-lag serial-correlation correction problem. Rather, it is a renewal-count
overdispersion problem: over a finite horizon, one symbol stream can realize many
short gaps while another realizes one or more very long gaps, even when the
nominal rates are equal. The resulting marginal imbalance can be much larger
than iid multinomial theory expects. These results are not LRD estimates; they
are control diagnostics that identify when a generator's long-memory mechanism
interferes with a simple marginal target.

\subsection{Toward One-Step Markov Validation}

Marginals are only the first local control.  The same validation philosophy
extends to one-step Markov behavior, but direct matrix testing is more
data-hungry.  A \(K\times K\) transition matrix has \(K^2\) cells, and the
effective count in row \(i\) is proportional to the number of visits to
symbol \(a_i\).  For large or skewed alphabets, many rows can have too little
data for a useful cell-wise chi-squared comparison.

A practical validation shortcut is to test interpretable low-dimensional
contrasts of the transition matrix before attempting a full matrix test. Useful
summaries include the repeat probability
\(\Pr(X_t=X_{t-1})\), stationary-weighted row total variation, selected rows
for scientifically important symbols, grouped transitions after aggregating
symbols into coarser classes, and projections \(u^\top(P-\hat P)v\) for a small
set of contrast vectors \(u\) and \(v\). These diagnostics do not replace full
row-by-row validation for small \(K\) or very long runs, but they provide a
shorter path for detecting whether local persistence, cycling, or group-level
movement has the intended strength.

The validation plots are mainly comparative. PB1 is the simplest baseline:
the latent spectral construction is visible and rank quantization often leaves
a recognizable symbolic power-law range. PB2 generally preserves a useful
long-memory signal, but the argmax transform and marginal-offset calibration
make the symbolic diagnostics less transparent than PB1. PB3 separates the
latent driver from the emitted Markov chain; this is useful because the latent
diagnostics can remain convincing even when the symbol indicators are partly blind to regime persistence. PB4 illustrates a different
phenomenology: dependence is created by intermittent laminar episodes rather
than by a Gaussian spectral pole.

Among model-based methods, MB1b and MB1c are intended as the most direct
symbolic additive-memory baselines, with MB1a kept as the exact finite-history
reference. MB2 and MB3 are mechanism-oriented renewal/regime baselines and are
particularly useful when recurrence or burst durations are the quantity of
interest. The weaker MB4a and MB5 results are included because MB4a's normalized
finite-history excitation can have a nearly white symbol spectrum, and MB5's copy/mutation growth can concentrate dependence into
local duplication structure rather than a clean one-hot power law. MB4b is the
replacement candidate motivated by this failure mode, but its finite-sample
behavior still requires validation. These cases show why symbolic LRD generators
require empirical diagnostics in addition to parameter labels.


\begin{figure*}[!t]
\centering
\acfsubfig{PB1_SpectralFGN}{PB1}\hfill
\acfsubfig{PB2_LGCM}{PB2}\hfill
\acfsubfig{PB3_WaveletMarkov_Haar}{PB3 Haar}\\[0.4ex]
\acfsubfig{PB3_WaveletMarkov_Spectral}{PB3 spectral}\hfill
\acfsubfig{PB4_IntermittentMapSymbols}{PB4}
\caption{Production ACF diagnostics for property-based methods. Each subplot is
the centered one-hot symbolic ACF with finite-length and target power-law guides;
latent ACFs and spectral-density panels are generated by the validation scripts
and synced into the paper figure tree, but are omitted here to save space.}
\label{fig:validation-acf-property}
\end{figure*}

\begin{figure*}[!t]
\centering
\acfsubfig{MB1a_LAMP}{MB1a}\hfill
\acfsubfig{MB1b_DyadicLAMP}{MB1b}\hfill
\acfsubfig{MB1c_CalibratedAdditiveMarkov}{MB1c}\\[0.4ex]
\acfsubfig{MB2_OnOffMarkov}{MB2}\hfill
\acfsubfig{MB3_FSS}{MB3}\hfill
\acfsubfig{MB4a_HawkesSymbol}{MB4a}\\[0.4ex]
\acfsubfig{MB4b_SelfExcitingMass}{MB4b}\hfill
\acfsubfig{MB4c_LogitSelfExcitingMass}{MB4c}\hfill
\acfsubfig{MB5_DuplicationMutation}{MB5}
\caption{Production ACF diagnostics for model-based methods. The panels use the
same title-free production plot style as Fig.~\ref{fig:validation-acf-property},
with method names supplied by subcaptions rather than embedded titles.}
\label{fig:validation-acf-model}
\end{figure*}

\section{Benchmark Results and Discussion}

The benchmark section supports the same operational narrative: what does each
mechanism cost when it is used as a generator rather than as a mathematical
idea?  The measurements are generated with \pkg{BenchmarkTools.jl} in a
separate environment.  The extended scaling
mode uses sequence lengths from $10^2$ to $10^6$ and averages each timing over
multiple independently seeded syntheses.  The resulting histograms and scaling
plots are written under \code{benchmark/results/} and summarized in
\code{benchmark/RESULTS.md}.

\benchfig{relative_times_k8_n1000000}{Relative generation time for alphabet
size $K=8$ and sequence length $n=10^6$. Times are normalized to the fastest
method in this benchmark run.}{benchmark-relative-k8}
\benchfig{scaling_k8}{Generation-time scaling with sequence length for
alphabet size $K=8$ on log-log axes.}{benchmark-scaling-k8}

\begin{table*}[!t]
\caption{\pkg{BenchmarkTools.jl} generation timings for $K=8$ and $n=10^6$ on the
development machine. Values are per-synthesis averages within each
\pkg{BenchmarkTools.jl} trial; the benchmark script used ten syntheses per trial. These
numbers are machine-specific and should be interpreted primarily as relative
costs for the configured standard cases.}
\label{tab:benchmark-k8}
\centering
\begin{tabular}{lrrrr}
\hline
Method & Time (ms) & Relative & Trial allocations & Trial memory (MiB) \\
\hline
\code{MB2} \code{OnOffMarkov} & 13.301 & 1.00 & 480 & 76.33 \\
\code{MB5} \code{DuplicationMutation} & 22.293 & 1.68 & 90 & 153.35 \\
\code{MB3} \code{FSS} & 23.347 & 1.76 & 60 & 76.30 \\
\code{PB4} \code{IntermittentMapSymbols} & 102.039 & 7.67 & 320 & 267.04 \\
\code{PB1} \code{SpectralFGN} & 117.242 & 8.81 & 490 & 648.52 \\
\code{PB3} \code{WaveletMarkov} spectral & 150.058 & 11.28 & 730 & 724.83 \\
\code{MB1c} \code{CalibratedAdditiveMarkov} & 270.006 & 20.30 & 80 & 152.59 \\
\code{MB4a} \code{HawkesSymbol} & 280.897 & 21.12 & 80 & 152.59 \\
\code{MB1a} \code{LAMP} & 283.760 & 21.33 & 80 & 152.59 \\
\code{MB4b} \code{SelfExcitingMass} & 284.050 & 21.36 & 80 & 152.59 \\
\code{PB2} \code{LGCM} & 439.012 & 33.01 & 1630 & 4348.82 \\
\code{MB4c} \code{LogitSelfExcitingMass} & 1390.853 & 104.57 & 100 & 152.59 \\
\code{MB1b} \code{DyadicLAMP} & 1557.556 & 117.10 & 130 & 762.98 \\
\hline
\end{tabular}
\end{table*}

The expected complexity classes are visible, although implementation constants
are also important.
Spectral methods pay for FFTs and sorting and scale near $O(n\log n)$.
Finite-history additive and self-exciting methods pay for their memory window.
Dyadic buckets trade exact history access for a scalable approximation, but the
current MB1b implementation still has a large constant factor.  The table
should therefore be read as engineering evidence, not as an asymptotic result.

The benchmark also reinforces the validation story.  MB2 and MB3 are fast
because they avoid long history scans, but their interpretation is tied to
renewal and regime mechanisms.  PB1 is a strong baseline because the latent
second-order target is simple, even though it pays for FFTs and allocation.  PB2
becomes expensive for large alphabets because it generates one latent stream per
symbol and calibrates offsets.  Speed only matters when the generator expresses
the intended symbolic dependence, and a convincing finite-sample plot is not a
proof of asymptotic LRD.
Both pieces of evidence are needed.


\section{Limitations and Future Work}

The main limitation is also the main reason for the package design: symbolic LRD
has several legitimate meanings.  Centered one-hot ACFs and spectra are transparent and
close to the DNA-indicator literature, but they can understate regime-based,
recurrence-based, or higher-order dependence.  Future validation should add
recurrence-time, count-variance, block-entropy, and local-structure diagnostics
without turning \sfive{} itself into an estimator package.

Several generators are intentionally research prototypes.  MB4a is retained as
a cautionary normalized-excitation baseline rather than a recommended generator.
MB4b and MB4c are stronger candidates, but still need a cleaner route from
power-law excitation to a stable symbolic spectral slope, perhaps through a
branching-process or count-process formulation before final categorical
sampling.  MB5 needs a better separation between local duplication patches and
long-lag copying.  MB3 needs a marginal-stabilized renewal construction.  A
promising variant would retain heavy-tailed renewal or intensity dynamics but
add finite-sample marginal calibration, for example by generating candidate
renewal events and then thinning, accepting, or allocating emissions against
target symbol counts.  That should probably be exposed as a calibrated MB3
variant rather than silently replacing the current FSS model, because the
current model's scientific identity is the competition among independent
fractal renewal streams.  PB3 would benefit from calibration that maximizes
contrast among regime stationary distributions while preserving requested local
transition matrices.

This is why the package is designed to expose weaknesses rather than hide them.
A symbolic LRD method should fail visibly when its mechanism does not support the
claim being made.  The next iteration should make those failures more diagnostic
and, where possible, turn them into better generators.

\section{Reproducibility Artifacts}

The paper is built from retained scripts and result artifacts.  Validation
figures are produced by scripts under \code{validation/}, with aggregate tables and plots
written under \code{validation/results/lrd\_diagnostics/}.  Benchmark tables and
plots are produced by \code{benchmark/benchmarks.jl} in its separate
\pkg{BenchmarkTools.jl} project, with extended scaling controlled by explicit
environment flags so rare runs are not part of ordinary tests.  The paper and
annotated bibliography are built by \code{Paper\_s4/Makefile}, which writes
LaTeX intermediates to \code{Paper\_s4/Tex/}.  The makefile also rsyncs the
production validation and benchmark plots from the package repository into
\code{Paper\_s4/Plots/} before compiling.

Reference provenance is part of the artifact.  Downloaded papers are kept
outside the public repository, while \code{Paper\_s4/DOWNLOAD.mn} records match
status and private-store paths.  The BibTeX database was checked with
\pkg{PaperFetch.jl} in online check-only mode; the check was used to correct
clear metadata conflicts without downloading additional PDFs.  This keeps the
public repository lightweight while preserving an audit trail for the literature
used to construct and interpret the generators.

\section{Conclusion}

\sfive{} is a reproducible workbench for synthesizing long-memory symbolic
data without treating symbols as numbers.  The package provides a common API,
explicit mechanisms, validation scripts, and benchmarks for property-based and
model-based generators.  It is useful for LRD
estimator testing, entropy and string-method test sets, controlled
machine-learning stress tests, synthetic event-log generation,
privacy-preserving simulation, anomaly and fake-source detection studies, and
workload modeling for symbolic streams.

The methodological point, and the narrative thread of the paper, is that for
symbolic LRD a parameter label is not enough.  We need to know the marginal, the
local structure, the mechanism, the transformation, the finite memory limit, and
the diagnostic used to judge the output.  Those pieces need to be explicit; no
single generator supplies all of them.  Future work should strengthen
correlation-calibrated symbolic generators, add exact numerical latent sources,
and connect the validation evidence to external estimator packages while keeping
\sfive{} focused on reproducible synthesis.
\appendices
\section{Slow Convergence of Averages Under LRD}

Long-range dependence changes how empirical averages stabilize. For an iid or
short-range dependent scalar process, the variance of a sample mean typically
scales as \(O(n^{-1})\), giving the familiar \(n^{-1/2}\) standard-error
rate. If a stationary scalar process \(Z_t\) has autocovariance
\begin{equation}
  \gamma(h) \sim c_\gamma h^{-\beta}, \qquad 0<\beta<1,
\end{equation}
then the covariance terms in the sample mean do not sum to a constant. Instead,
\begin{align}
  \operatorname{Var}(\bar Z_n)
  &= \frac{1}{n^2}
     \left[n\gamma(0)+2\sum_{h=1}^{n-1}(n-h)\gamma(h)\right] \\
  &\sim C n^{-\beta},
\end{align}
for a positive constant \(C\) when the tail coefficient is positive
\cite{beran1994_statistics,samorodnitsky2006_long,pipiras2017_long}. Thus the
standard error decays as \(n^{-\beta/2}\), rather than \(n^{-1/2}\). In
Hurst-parameter notation, \(\beta=2-2H\), so the standard error behaves like
\(n^{H-1}\). For \(H=0.8\), for example, the standard error is of order
\(n^{-0.2}\), so increasing the sequence length by a factor of 100 only reduces
the typical average error by about a factor of \(100^{-0.2}\approx 0.40\),
not by a factor of 10.

The same calculation applies directly to marginal estimation for symbolic
sequences. For alphabet symbol \(a_j\), define the indicator
\begin{equation}
  I_{j,t}=\mathbf{1}\{X_t=a_j\}, \qquad
  \hat p_j = \frac{1}{n}\sum_{t=1}^{n} I_{j,t}.
\end{equation}
If the centered indicator process \(I_{j,t}-p_j\) has autocovariance
\(\gamma_j(h)\sim c_jh^{-\beta_j}\), then
\begin{equation}
  \operatorname{Var}(\hat p_j) \sim C_j n^{-\beta_j}.
\end{equation}
This is much larger than the iid multinomial approximation
\(p_j(1-p_j)/n\) at large \(n\). Equivalently, the effective number of
independent observations grows sublinearly; up to constants, an LRD sequence
with exponent \(\beta\) can behave like only \(n_{\mathrm{eff}}\propto
n^\beta\) independent observations for mean estimation. The exact constant
depends on the full autocorrelation function, not only on its exponent, so this
scaling should be read as an order-of-magnitude guide rather than a replacement
for a generator-specific uncertainty calculation.

This has two practical consequences for the marginal-control validation in
Section~\ref{sec:validation}. First, a histogram from a long symbolic sequence
can still show visible departures from the requested marginal even when the
generator has the right long-run target. These departures need not disappear at
the iid \(n^{-1/2}\) rate. Second, a classical multinomial chi-squared test
treats the counts as though they came from \(n\) independent draws. For LRD
sequences this can make p-values too small and can overstate the evidence
against the target marginal. The ESS-adjusted statistic used here is an approximate
correction: it reduces the nominal sample size using the observed one-hot
autocorrelation. It remains approximate, because a full dependent multinomial
test would require the joint covariance structure of the count vector, but it is
a guard against interpreting slow LRD convergence as an immediate
marginal-control failure.

Uniform targets are still valuable in this setting. They keep expected bin
counts away from zero and make deviations visually comparable across symbols.
However, the right interpretation is diagnostic rather than absolute: rank-binned
methods may match the full-sample marginal almost exactly while the trimmed
interior fluctuates, and process-based methods may require many independent
replicates before their average histogram settles. Trimming the first and last
10\% reduces startup and terminal edge effects, but it does not remove the slow
convergence caused by long memory itself.


\subsection{Uniform Marginal Histograms}

The validation script writes a complete histogram grid for all methods, but the
paper includes only the worst trimmed-total-variation case from the current
validation run. The remaining methods are summarized by the marginal table in
the main validation section and have smaller deviations from the target uniform
marginal. Figure~\ref{fig:uniform-marginal-worst} shows the mean trimmed
empirical symbol frequencies for that worst case; error bars show one replicate
standard deviation and the dashed red line marks the intended frequency \(1/K\).

\begin{figure*}[!t]
\centering
\IfFileExists{\marginalpath/production/marginal/MB3_FSS_uniform_marginal_histogram_k8.pdf}{%
  \includegraphics[width=0.72\textwidth]{\marginalpath/production/marginal/MB3_FSS_uniform_marginal_histogram_k8.pdf}%
}{%
  \fbox{\parbox{0.72\textwidth}{\centering Missing generated marginal production
  histogram PDF. Run \code{julia --project=. validation/marginal\_control.jl}.}}%
}
\caption{Worst-case uniform marginal validation histogram for \(K=8\) in the
current validation run. The full generated grid is retained as a validation
artifact, but only this panel is included here because the other methods are
closer to the target marginal.}
\label{fig:uniform-marginal-worst}
\end{figure*}

\section{Transformations of LRD Processes}

Let $Z_t$ be a real-valued, zero-mean stationary process with autocovariance
\begin{equation}
  \gamma_Z(k) \sim c k^{-\beta}, \qquad 0 < \beta < 1.
\end{equation}
For a linear transformation $Y_t = aZ_t + b$ with $a \ne 0$,
\begin{equation}
  \gamma_Y(k) = a^2 \gamma_Z(k),
\end{equation}
so the LRD exponent is unchanged.

\subsection{Linear Maps and Vector Encodings}

For vector-valued encodings, let \(\vect{Y}_t = B\vect{e}(X_t)\), where
\(\vect{e}(a_i)\) is the
\(i\)th coordinate vector and \(B\) is a fixed matrix. If the centered indicator
covariance matrix
\begin{equation}
  \Gamma_I(k)=
  \mathbb{E}\left[
    (\vect{e}(X_t)-\vect{p})(\vect{e}(X_{t+k})-\vect{p})^\top
  \right]
\end{equation}
has entries with a common power-law tail, then
\begin{equation}
  \Gamma_{\vect{Y}}(k)=B\Gamma_I(k)B^\top.
\end{equation}
A scalar numerical encoding is the special case \(B=\vect{b}^\top\). This can preserve a
tail if \(\vect{b}\) has nonzero projection on the long-memory eigenspace of
\(\Gamma_I(k)\), but it can also cancel or hide dependence. This is one reason
arbitrary scalar DNA walks or word-index encodings can be misleading, and why
\sfive{} defaults to centered one-hot diagnostics \cite{voss1992_evolution}.

\subsection{Gaussian Nonlinear Transformations}

For nonlinear transformations, preservation depends on both the function and the
law of the latent process. The cleanest classical result is for a stationary
standard Gaussian process $Z_t$ with correlation
\begin{equation}
  r(k) \sim c k^{-\beta}.
\end{equation}
Let $h$ be square-integrable under the standard normal law and centered so that
\(\mathbb{E}h(Z)=0\). Its Hermite expansion is
\begin{equation}
  h(x) = \sum_{q=0}^{\infty} \frac{J_q}{q!} H_q(x),
  \qquad
  J_q=\mathbb{E}[h(Z)H_q(Z)],
\end{equation}
where $H_q$ is the $q$th Hermite polynomial. The Hermite rank $m$ is the first
positive $q$ with $J_q\ne0$. Using orthogonality of Hermite polynomials for
joint Gaussian variables,
\begin{equation}
  \operatorname{Cov}(h(Z_t),h(Z_{t+k}))
  =
  \sum_{q=m}^{\infty} \frac{J_q^2}{q!} r(k)^q.
\end{equation}
Thus, when the first nonzero term dominates,
\begin{equation}
  \operatorname{Cov}(h(Z_t),h(Z_{t+k}))
  \sim
  \frac{J_m^2}{m!} c^m k^{-m\beta}.
\end{equation}
The transformed process remains covariance-LRD only when $m\beta<1$.
Otherwise the covariance tail becomes summable even though the latent process
was LRD. This Hermite-rank principle underlies classical non-central limit
theory for nonlinear functions of Gaussian LRD processes
\cite{dobrushin1979_non,taqqu1975_weak,breuer1983_central,arcones1994_limit}.

\subsection{Thresholds, Quantization, and Indicators}

Quantization is a nonlinear map. For a binary threshold
\begin{equation}
  Y_t=\mathbf{1}\{Z_t\le q\}-\Phi(q),
\end{equation}
the first Hermite coefficient is
\begin{equation}
  J_1 = \mathbb{E}[Y_t Z_t] = -\phi(q),
\end{equation}
which is nonzero for finite $q$. Hence a single Gaussian threshold usually has
Hermite rank one and preserves the covariance exponent:
\begin{equation}
  \operatorname{Cov}(Y_t,Y_{t+k}) \sim \phi(q)^2 r(k).
\end{equation}
For a multi-bin quantizer, the centered indicator for bin
\((q_{i-1},q_i]\) has first Hermite coefficient
\begin{equation}
  J_{1,i} =
  \mathbb{E}\left[
      \left(\mathbf{1}\{q_{i-1}<Z\le q_i\}-p_i\right)Z
  \right]
  =
  \phi(q_{i-1})-\phi(q_i),
\end{equation}
with the convention that $\phi(\pm\infty)=0$. Interior symmetric bins can have
small or zero first coefficient, in which case higher Hermite terms can dominate
at observable scales. Rare bins can also have weak finite-sample diagnostics
because their indicator variance is small and long-lag estimates have high
relative noise.

This analysis explains why PB1 can often preserve a visible power-law in common
symbol indicators, while some symbols or transformations may look weaker. It
also explains why PB2's argmax transform is harder to reason about: the emitted
indicator is a nonlinear function of a $K$-dimensional Gaussian vector, and its
effective Hermite rank depends on the offsets, marginal probabilities, and
symmetries among latent streams.

\subsection{Markov Regime Transformations}

PB3 introduces another layer: a latent LRD driver determines a regime process
$R_t$, and symbols are emitted through regime-specific Markov chains. If
\(R_t\) has long memory and the stationary distribution vectors
\(\vect{\pi}^{(r)}\) of the
matrices $P^{(r)}$ differ across regimes, then a symbol indicator can inherit
part of the regime dependence through
\begin{equation}
  \mathbb{E}[\mathbf{1}\{X_t=a_j\}\mid R_t=r]\approx \pi^{(r)}_j.
\end{equation}
However, if all regimes have the same stationary distribution, the first-order
symbol indicators may be nearly blind to the regime process even though
higher-order local structure is changing. In that case, diagnostics based only
on one-hot autocorrelation can understate the dependence created by the latent
driver. This motivates validating both latent regimes and emitted symbols.

\subsection{Non-Gaussian Latent Processes}

The Hermite argument does not apply directly to intermittent maps or
heavy-tailed renewal mechanisms. For such processes, transformations can alter
dependence through changes in laminar-state visibility, tail truncation, and
finite-sample scale range. Intermittent-map quantization may emphasize laminar
episodes if a bin isolates the sticky region near zero, or weaken them if rank
binning spreads laminar values across symbols. Heavy-tailed on/off and renewal
models can preserve count-process LRD while individual symbol autocorrelations
are noisy or dominated by rare-event effects.

The practical conclusion for \sfive{} is conservative. The package
does not claim that a nominal parameter always determines the finite symbolic
ACF exponent. Instead it records the mechanism, exposes latent outputs where
they exist, and produces validation evidence at the numerical and symbolic
levels.

\bibliographystyle{IEEEtran}
\bibliography{references/s5_references}

\end{document}
